\begin{table}[ht]
\centering
\caption{Nombre d’essais inclus dans les contrastes de confiance à percept apparié.}
\label{tab:matched_confidence_trial_counts}
\footnotesize
\begin{tabular}{llcc}
\toprule
Analyse & SOA & Participants congr./illus. & Essais congr./illus. \\
\midrule
Fission : percept 2 flashs & Standard & 9 / 9 & 637 / 229 \\
Fission : percept 2 flashs & Court & 9 / 9 & 541 / 315 \\
Fusion : percept 1 flash & Standard & 9 / 8 & 1214 / 107 \\
Fusion : percept 1 flash & Court & 9 / 9 & 1214 / 338 \\
\bottomrule
\end{tabular}

\vspace{0.2cm}
\begin{minipage}{0.92\textwidth}
\footnotesize
\textit{Note.} Les colonnes indiquent le nombre de participants et le nombre d’essais contribuant respectivement à la condition congruente puis à la condition illusoire. Pour la fusion, la condition congruente A1V1 ne comporte pas de SOA et sert donc de référence commune aux contrastes avec les conditions de fusion standard et courte.
\end{minipage}
\end{table}
